% RESEARCH_DRAFT_SHA256: 645793364d7bc2e09d7b56f8e9abf50659c728e574899e1327959b2965f617da
% Layout and source-reference wrappers expanded for the web renderer.
\documentclass[11pt]{article}
\usepackage[margin=25mm]{geometry}
\usepackage{fontspec,unicode-math}
\setmainfont{Latin Modern Roman}
\setmathfont{Latin Modern Math}
\usepackage{amsmath,amssymb,mathtools,xurl,hyperref}
\hypersetup{colorlinks=true,urlcolor=blue,linkcolor=blue}
\setlength{\parindent}{0pt}
\setlength{\parskip}{5pt}
\setlength{\emergencystretch}{3em}
\newcommand{\sref}[2]{\hyperlink{source:#1:#2}{[S#2]}}
\newcommand{\cataloguescope}[1]{\paragraph{Source question.}#1}
\newcommand{\floor}[1]{\left\lfloor#1\right\rfloor}
\begin{document}

\section*{Constructive frame partitioning: commuting inputs and a dimension-dependent barrier}
\noindent UnsolvedMath ID 2800605; AMR-027-0605. Research attempt dated 11 October 2026.
\paragraph{Result in scope.}This manuscript contains a source-specific attempt and complete proofs of the stated auxiliary calculations. It makes no new solution or novelty claim for the entire target.
\section{Original problem and preserved definition}
The following definition and source list are reproduced from the repository. Any earlier inspection dates and reported literature claims in this preserved text describe that source record, not additional verification by this attempt.
% BEGIN PRESERVED DEFINITION
\subsection{Definitions and mathematical statement}

\paragraph{Shared notation.}$\mathbb N=\{1,2,\ldots\}$, and $\mathbb Z,\mathbb Q,\mathbb R,\mathbb C$ denote the integers, rationals, reals and complex numbers. Any other conventions are specified locally.


A tight frame at level \(\eta\) is a finite family \(v_1,\ldots,v_m\in\mathbb C^d\) satisfying
\[\sum_{i=1}^m v_i v_i^*=\eta I_d,\]
equivalently \(\sum_i|\langle u,v_i\rangle|^2=\eta\) for every unit vector \(u\in\mathbb C^d\). Here \(v_i^*\) denotes conjugate transpose and \(\|\cdot\|_2\) is the Euclidean norm. Weaver's two-partition conclusion asks for a partition \([m]=S_1\sqcup S_2\) satisfying
\[\sum_{i\in S_j}|\langle u,v_i\rangle|^2\le\eta-\theta
\quad\text{for every unit }u\text{ and }j=1,2,\]
with universal constants \(\eta\ge2\) and \(\theta>0\), when \(\|v_i\|_2\le1\). Marcus, Spielman and Srivastava established the existence of such partitions; for example their result allows \(\eta=18\) and \(\theta=2\).
\paragraph{Statement recorded in the supplied source.}
Give a polynomial-time construction of a partition with Weaver's two-partition property: find universal constants \(\eta\ge2\), \(\theta>0\) and an algorithm that, for every finite tight frame as above with \(\|v_i\|_2\le1\), outputs \(S_1,S_2\) satisfying the displayed inequality in time polynomial in the size of its input. In particular, can the partitions guaranteed by the Kadison--Singer existence theorem be found efficiently? \href{https://arxiv.org/pdf/2205.02161}{[S2]}
\subsection{Short English statement}
Construct efficiently the two-partition of any bounded tight frame guaranteed by Weaver’s formulation of the Kadison–Singer theorem, using universal constants.
\subsection{Sources}
\begin{description}
\item[S1] ULAM.AI and the UnsolvedMath Contributors, UnsolvedMath: A Curated Collection of Open Mathematics Problems, record 2800605 (AMR-027-0605). CC BY 4.0. \url{https://huggingface.co/datasets/ulamai/UnsolvedMath}.
\item[S2] Ben Jourdan, Peter Macgregor and He Sun, Is the Algorithmic Kadison-Singer Problem Hard? (2022), arXiv:2205.02161, Section 1, Conjecture 1 and the algorithmic question immediately following it. \url{https://arxiv.org/pdf/2205.02161}.
\end{description}
% END PRESERVED DEFINITION
\subsection{Verbatim supplied catalogue question}
Give a (polynomial time) construction of the tight frame partition satisfying the properties required in the Kadison-Singer problem (or the related Weaver’s conjecture). These partitions were proven to exist (with a non-constructive proof) in the recent breakthrough of Marcus, Spielman, and Srivastava .

\section{Attempt and results}
Keep $\sum_i v_iv_i^*=\eta I_d$ and $\|v_i\|\le1$. An algorithm for the
full target must use constants $\eta,\theta$ independent of both $d$ and $m$.
The following algorithms use exact real/complex arithmetic; polynomial bit
complexity for approximate inputs is not claimed.

\paragraph{A constructive commuting special case.}
Suppose the rank-one matrices $v_iv_i^*$ commute pairwise. Two nonzero such
matrices either have orthogonal ranges or the same range. To see this, write
them $vv^*$ and $ww^*$. If $v^*w=0$ they multiply to zero in either order.
Otherwise their equal nonzero products have ranges respectively $\mathbb Cv$
and $\mathbb Cw$, forcing these ranges to agree. Thus the vectors split
into groups along mutually orthogonal directions. Tightness implies that
these directions span the entire space and that the sum of squared lengths
in each direction is $\eta$.
In each group, process weights $a_i=\|v_i\|^2\in[0,1]$, assigning the next
weight to the smaller of two current loads. If their difference is $D\le1$,
the next difference is $|D-a_i|\le1$. Starting at zero proves that the final
loads differ by at most one, hence each is at most $(\eta+1)/2$.
Grouping parallel/orthogonal vectors by inner products and applying this rule
takes polynomially many arithmetic operations. Zero vectors may be placed
arbitrarily. For $\eta=2$ this gives the source conclusion with $\theta=1/2$
for every commuting input, in every dimension. The proof covers all unit test
vectors because the two partial frame operators are diagonal in the
orthonormal direction basis and each diagonal entry is at most $3/2$.

\paragraph{Attempt to extend by a matrix-energy potential.}
Set $A_i=v_iv_i^*$. We construct signs $\epsilon_i\in\{-1,1\}$ so that
$S=\sum_i\epsilon_iA_i$ has small Frobenius norm. At a partial sum $T$,
\[
\tfrac12\|T+A_i\|_F^2+\tfrac12\|T-A_i\|_F^2
=\|T\|_F^2+\|A_i\|_F^2.
\]
Choose the sign attaining the smaller value. Induction yields
\[
\|S\|_{\mathrm{op}}^2\le\|S\|_F^2
\le\sum_i\|A_i\|_F^2=\sum_i\|v_i\|^4
\le\sum_i\|v_i\|^2=\eta d.
\]
The two partition operators are $(\eta I\pm S)/2$, and thus have norms
at most $(\eta+\sqrt{\eta d})/2$. This is a deterministic polynomial
arithmetic algorithm, but its guarantee implies a positive saving only when
$\eta>d$. Choosing $\eta=4d$ gives a saving at least $\eta/4$, yet violates
the original requirement of universal constants. For fixed $\eta$ and
unbounded $d$ the guarantee eventually becomes vacuous. This pinpoints the
failure of the attempted Frobenius-norm derandomization.

\section{Context and adversarial audit}
The primary abstract of Jourdan, Macgregor and Sun, arXiv:2205.02161, was
consulted on 11 October 2026. Its algorithmic and hardness conclusions have
specific parameter regimes; its stated FNP-hard parameter is not used to
rule out the source's universal-constant existence construction. The
commuting and energy calculations here are elementary special-case/known
methods; no new general algorithm or novelty is claimed. The audit checked
zero vectors, complex conjugation, operator versus Frobenius norm, and the
dimension-independence requirement. No numerical or formal computation was
performed for this record.

\section{Process and contributions}
Mathematics and model-assisted adversarial self-audit: OpenAI gpt-6-astra, ultra reasoning, as dispatched for this research batch; exact serving revision not exposed. Codex (GPT-6; exact serving revision not exposed) materially assisted with orchestration and format review. No human mathematical authorship, mathematical edits, or independent human proof checking is asserted.
Contribution prompt A was applied to the full definition and supplied source row, then prompt B was applied adversarially to the derived results. The item-specific audit above records the scope checks and rejected extrapolations. Local assembly and structural checks do not establish mathematical correctness. The source catalogue status is preserved in TOP\_PROBLEM and is separate from this attempt status.
\paragraph{Integration review.} Codex (GPT-6; exact serving revision not exposed) reviewed the written arguments and expanded layout and source-reference wrappers. Repository metadata and the production web renderer were checked using Python 3.13 and JavaScript with MathJax 3.2.2. This model-assisted review is not independent human proof verification; no PDF compilation or formal proof checking is asserted.

\section{Completion Estimate}
\noindent\textbf{COMPLETION: 0\%}
The proved material is a rederivation of known or auxiliary facts rather than a substantiated new advance on the original target. The energy proof loses sqrt(d), and the commuting reduction excludes general noncommuting tight frames; no algorithm with universal constants for all inputs is obtained.
\section{References}
The complete original Sources subsection above is retained, including catalogue attribution and licence notices. The additional list records primary sources actually consulted for this attempt on 11 October 2026; uncited source links reproduced above were not all rechecked.
\begin{itemize}
\item Ben Jourdan, Peter Macgregor and He Sun, Is the Algorithmic Kadison-Singer Problem Hard?, arXiv:2205.02161; abstract and parameter qualifications inspected. \url{https://arxiv.org/abs/2205.02161}.
\end{itemize}
\section{Final status}
\textbf{Unresolved in this attempt.} The energy proof loses sqrt(d), and the commuting reduction excludes general noncommuting tight frames; no algorithm with universal constants for all inputs is obtained.
\end{document}
